% Include-only section; requires amsmath, amssymb, booktabs.
\section{A control-flow error in the archived scratch-biased experiment}
\label{sec:scratch-correction}

The scratch-biased 215-state experiment in
\texttt{T215\_Transfer\_2026-10-20.zip} has a successful quantified
primitive-runtime certificate, but its README's high-level command
description is incorrect. Replacing a continuation leaf by a two-leaf
subtree changed the continuation's target. This defect is visible in
the literal source independently of physical allocation.

The pre-insertion continuation was $\mathsf{CC}(3)$ at binary path
$10011$. Its three right ancestors are at indices $0,3,4$, so the
third outward right ancestor is the root and the continuation returns
to source ordinal zero. The replacement subtree
$(\mathsf{CI}(6),\mathsf{CC}(3))$ places the continuation at $100111$.
There are now four right ancestors, at indices $0,3,4,5$. Its third
outward right ancestor has prefix $100$, whose first leaf is source
ordinal $23$. Thus the literal edge is
\[
       4697:\ \mathsf{CC}(3)\longrightarrow23,
\]
not an edge to zero. The added right branch must be counted in a
relative continuation instruction.

\paragraph{Exact unconditional halting suffix.}
At ordinal $4696$, assume only that logical registers
$A=H=X=0$; write the arbitrary current top as $T=t$. Other logical
registers are irrelevant to the following suffix. The actual execution
is:
\begin{enumerate}
 \item Ordinal $4696$ increments scratch $X$ to one; ordinal $4697$
 returns to $23$.
 \item The zero selector $A$ takes the zero branch at $23$, reaching
 ordinal $4698$, the top-clear loop. That loop clears $T$ in exactly
 $2t+1$ source steps and does not change $H$ or $X$.
 \item The copy-cache code begins at $4700$. Its $H$ loop is skipped
 because $H=0$, while its scratch-restoration loop consumes the
 outstanding $X=1$ and increments $H$ to one.
 \item The inherited NIP test begins at $4707$ and recognizes
 precisely this $H=1$. It reaches the halt primitive at ordinal
 $4718$.
\end{enumerate}
Counting transitions from $4696$ to the halt primitive gives exactly
\[
                            2t+15.
\]
The proof is a single unary-loop induction plus the eleven fixed
post-clear edges. No inference about a very long simulated prefix is
needed.

Both AX and MP return to $4696$ with $A=H=X=0$ according to their
unchanged arithmetic contracts. In particular, the supplied 69,744-step
blank trace and 172-step AX-prefix match put the exact table inside
the AX begun with registers $[1,0,3,0,0,0,0,0,0]$. Completing that AX
returns with gate/top zero and the required zero selector/cache/scratch.
The source then reaches its halt primitive after fifteen further source
transitions. Subject to the supplied complete primitive simulation, the
archived exact binary table therefore halts on blank tape
\emph{unconditionally}. It cannot provide the requested nonhalting
equivalence when ZFC is consistent.

The correction does not dispute the primitive runtime checker: that
checker correctly implements the literal source. The failed boundary
is between the literal source and the advertised high-level parser.
Consequently the proposed first-field bias, candidate steps of one or
two, and cofinality problem do not describe this exact machine. Those
ideas may remain relevant to a repaired source that actually returns
to its scheduler, but cannot calibrate these bytes.

The accompanying independent checker resolves both complete guarded
control graphs, verifies the changed edge, records the exact fifteen-step
suffix at $t=0$, and checks the affine count for $0\le t\le31$.
The finite count check corroborates the explicit unary-loop proof; it
does not replace that proof or rerun the full physical certificates.
